\chapter{Les fonctions en C}

\begin{synthese}[Au programme]
Découper un programme en \textbf{fonctions} : déclarer un prototype, écrire une
définition, appeler la fonction et lui passer des \textbf{paramètres}. Comprendre
le rôle de \code{return}, la différence entre variables \textbf{locales} et
\textbf{globales}, et maîtriser la \textbf{récursivité} (factorielle, puissance,
Fibonacci). C'est l'outil qui rend un programme C lisible et réutilisable.
\end{synthese}

\section{Pourquoi des fonctions ?}

Jusqu'ici, tout notre code tenait dans la fonction \code{main}. Mais dès qu'un
programme grandit, écrire tout d'un seul bloc devient illisible et on recopie
sans cesse les mêmes calculs. La solution : isoler chaque traitement dans une
\textbf{fonction} qu'on appelle au besoin.

\begin{definition}
Une \textbf{fonction} est un sous-programme nommé qui réalise une tâche précise.
On l'écrit \textbf{une seule fois}, puis on l'\textbf{appelle} autant de fois que
nécessaire. Elle peut recevoir des données (les \textbf{paramètres}) et renvoyer
un résultat (la valeur de \textbf{retour}).
\end{definition}

Trois bénéfices immédiats :
\begin{itemize}
  \item \textbf{Réutiliser} : on écrit le calcul de la moyenne une fois, on s'en
        sert pour chaque élève du lycée.
  \item \textbf{Structurer} : un gros problème se découpe en petites fonctions
        simples, faciles à lire et à corriger.
  \item \textbf{Tester} : chaque fonction se vérifie séparément.
\end{itemize}

\begin{remarque}
Tu connais déjà des fonctions : \code{printf} et \code{scanf} de la bibliothèque
\code{<stdio.h>}, ou \code{sqrt} de \code{<math.h>}, sont des fonctions écrites
par d'autres et que tu appelles. Tu vas maintenant écrire les tiennes.
\end{remarque}

\section{Déclarer une fonction}

\subsection{La définition}

Définir une fonction, c'est écrire son code complet. La structure est toujours
la même.

\begin{syntaxe}
\code{type nom(parametres)} \code{\{}\ \textit{instructions}\ \code{return valeur;}\ \code{\}}
\end{syntaxe}

\begin{itemize}
  \item \code{type} : le type de la valeur renvoyée (\code{int}, \code{float},
        \code{double}, \code{char}\dots ou \code{void} si la fonction ne renvoie rien) ;
  \item \code{nom} : le nom de la fonction (mêmes règles qu'une variable) ;
  \item \code{parametres} : la liste des données reçues, séparées par des virgules,
        chacune précédée de son type ;
  \item entre accolades : le \textbf{corps} de la fonction ;
  \item \code{return} : renvoie le résultat \textbf{et} termine la fonction.
\end{itemize}

\begin{codec}
// Fonction qui calcule le carre d'un entier
int carre(int x) {
    int resultat = x * x;   // on calcule le carre
    return resultat;        // on renvoie le resultat
}
\end{codec}
\fig{Ici \texttt{int} est le type de retour, \texttt{carre} le nom, \texttt{int x} le paramètre.}

\subsection{Le rôle de \texttt{return}}

\begin{definition}
L'instruction \code{return} fait deux choses : elle \textbf{renvoie} une valeur à
l'endroit qui a appelé la fonction, et elle \textbf{arrête} immédiatement la
fonction (les instructions situées après ne sont pas exécutées).
\end{definition}

La valeur renvoyée par \code{return} doit être du même type que celui annoncé
devant le nom de la fonction. Une fonction \code{int} renvoie un entier, une
fonction \code{float} renvoie un décimal.

\begin{exemple}
Dans \code{int carre(int x)}, l'appel \code{carre(5)} vaut \code{25}. On peut
l'utiliser comme une valeur : \code{int s = carre(5) + carre(3);} place \code{34}
dans \code{s}.
\end{exemple}

\subsection{Les fonctions \texttt{void} (procédures)}

Certaines fonctions n'ont pas de résultat à renvoyer : elles se contentent
d'\emph{agir} (afficher quelque chose, par exemple). Leur type de retour est
alors \code{void} (« vide »), et elles n'utilisent pas \code{return valeur;}.
On les appelle souvent des \textbf{procédures}.

\begin{codec}
// Procedure : elle affiche, mais ne renvoie rien
void afficherLigne(void) {
    printf("--------------------\n");
}
\end{codec}

\begin{remarque}
Le mot \code{void} entre les parenthèses indique que la fonction ne prend
\textbf{aucun} paramètre. On peut aussi laisser les parenthèses vides :
\code{void afficherLigne()}.
\end{remarque}

\section{Appeler une fonction}

\textbf{Appeler} une fonction, c'est l'exécuter en écrivant son nom suivi des
valeurs à lui transmettre entre parenthèses (les \textbf{arguments}).

\begin{codec}
#include <stdio.h>

int carre(int x) {
    return x * x;
}

int main(void) {
    int n = 6;
    int c = carre(n);              // appel : on passe 6, on recupere 36
    printf("Le carre de %d est %d\n", n, c);
    printf("Le carre de 4 est %d\n", carre(4));  // appel direct
    return 0;
}
\end{codec}
\fig{Affiche : « Le carre de 6 est 36 » puis « Le carre de 4 est 16 ».}

\subsection{Le passage des paramètres par valeur}

\begin{propriete}[Passage par valeur]
En C, les paramètres sont passés \textbf{par valeur} : la fonction reçoit une
\textbf{copie} de l'argument. Modifier le paramètre à l'intérieur de la fonction
\textbf{ne change pas} la variable d'origine de l'appelant.
\end{propriete}

\begin{codec}
#include <stdio.h>

void doubler(int x) {
    x = x * 2;                 // on modifie la COPIE locale
    printf("Dans la fonction : x = %d\n", x);
}

int main(void) {
    int a = 10;
    doubler(a);                // on passe une copie de a
    printf("Dans main : a = %d\n", a);   // a vaut toujours 10 !
    return 0;
}
\end{codec}
\fig{Affiche « Dans la fonction : x = 20 » puis « Dans main : a = 10 ».}

\begin{attention}
La variable \code{a} de \code{main} n'est pas modifiée : la fonction travaille sur
sa propre copie. Pour qu'une fonction modifie réellement une variable de
l'appelant, il faut les \textbf{pointeurs} (chapitre dédié) ; ce n'est pas l'objet
de ce chapitre.
\end{attention}

\subsection{Vocabulaire : paramètre ou argument ?}

\begin{itemize}
  \item le \textbf{paramètre} est la variable écrite dans la définition
        (\code{int x} dans \code{carre(int x)}) ;
  \item l'\textbf{argument} est la valeur réellement transmise lors de l'appel
        (\code{6} dans \code{carre(6)}).
\end{itemize}

\section{Où placer les fonctions : le prototype}

Le compilateur lit le fichier de haut en bas. Il doit \textbf{connaître} une
fonction avant qu'on l'appelle. Deux solutions.

\begin{enumerate}[label=\arabic*)]
  \item Écrire la définition complète \textbf{avant} \code{main}.
  \item Écrire seulement le \textbf{prototype} avant \code{main}, et la définition
        complète après. C'est la méthode recommandée car \code{main} reste en haut.
\end{enumerate}

\begin{definition}
Un \textbf{prototype} (ou \emph{déclaration}) annonce au compilateur le nom de la
fonction, son type de retour et le type de ses paramètres, suivi d'un
point-virgule. Il ne contient pas le corps.
\end{definition}

\begin{syntaxe}
\code{type nom(types des parametres);}\quad ex.\quad \code{int carre(int x);}
\end{syntaxe}

\begin{codec}
#include <stdio.h>

int carre(int x);        // PROTOTYPE (avant main)

int main(void) {
    printf("%d\n", carre(7));   // l'appel est valide grace au prototype
    return 0;
}

int carre(int x) {       // DEFINITION (apres main)
    return x * x;
}
\end{codec}

\begin{remarque}
Dans un prototype, on peut omettre le nom des paramètres :
\code{int carre(int);} est correct. On garde souvent les noms pour la lisibilité.
\end{remarque}

\section{Variables locales et globales : la portée}

\begin{definition}
La \textbf{portée} (\emph{scope}) d'une variable est la zone du programme où cette
variable existe et peut être utilisée.
\end{definition}

\subsection{Variables locales}

Une variable déclarée \textbf{à l'intérieur} d'une fonction est \textbf{locale} :
elle n'existe que pendant l'exécution de cette fonction et reste invisible aux
autres. Les paramètres sont également des variables locales.

\begin{codec}
#include <stdio.h>

void fonctionA(void) {
    int compteur = 5;    // locale a fonctionA
    printf("A : %d\n", compteur);
}

int main(void) {
    int compteur = 100;  // une AUTRE variable locale a main
    fonctionA();
    printf("main : %d\n", compteur);   // affiche 100
    return 0;
}
\end{codec}
\fig{Les deux \texttt{compteur} portent le même nom mais sont indépendants.}

\subsection{Variables globales}

Une variable déclarée \textbf{en dehors} de toute fonction (souvent en haut du
fichier) est \textbf{globale} : toutes les fonctions peuvent la lire et la
modifier.

\begin{codec}
#include <stdio.h>

int total = 0;           // variable GLOBALE

void ajouter(int montant) {
    total = total + montant;   // accede a la globale
}

int main(void) {
    ajouter(2000);       // en francs CFA
    ajouter(3500);
    printf("Total : %d FCFA\n", total);   // affiche 5500
    return 0;
}
\end{codec}

\begin{attention}
Les variables globales sont à utiliser \textbf{avec parcimonie} : comme n'importe
quelle fonction peut les modifier, elles rendent les erreurs difficiles à
retrouver. Préfère toujours des variables locales et le passage de paramètres.
\end{attention}

\section{Quelques fonctions utiles}

\subsection{Maximum de deux entiers}

\begin{codec}
// Renvoie le plus grand des deux entiers a et b
int max(int a, int b) {
    if (a > b) {
        return a;        // return arrete la fonction ici
    }
    return b;            // atteint seulement si a <= b
}
\end{codec}

\subsection{Tester la parité}

\begin{codec}
// Renvoie 1 si n est pair, 0 sinon
int estPair(int n) {
    return (n % 2 == 0);   // la condition vaut 1 (vrai) ou 0 (faux)
}
\end{codec}

\begin{remarque}
En C, une condition vraie vaut \code{1} et une condition fausse vaut \code{0}.
On peut donc directement renvoyer le résultat d'un test, comme ci-dessus.
\end{remarque}

\subsection{Moyenne de trois notes}

\begin{codec}
#include <stdio.h>

// Moyenne de trois notes (valeurs decimales)
float moyenne(float a, float b, float c) {
    return (a + b + c) / 3.0;
}

int main(void) {
    float m = moyenne(12.5, 9.0, 15.5);
    printf("Moyenne : %.2f / 20\n", m);   // affiche 12.33 / 20
    return 0;
}
\end{codec}
\fig{Le type de retour \texttt{float} permet de garder les décimales.}

\section{La récursivité}

\subsection{Le principe}

\begin{definition}
Une fonction est dite \textbf{récursive} lorsqu'elle s'\textbf{appelle elle-même}.
Pour qu'elle s'arrête, elle doit comporter deux parties :
\begin{itemize}
  \item un \textbf{cas de base} : une situation simple où le résultat est connu
        directement, sans nouvel appel ;
  \item un \textbf{cas récursif} : la fonction se rappelle sur un problème
        \textbf{plus petit}, qui se rapproche du cas de base.
\end{itemize}
\end{definition}

\begin{attention}
Sans cas de base, ou si les appels ne se rapprochent jamais du cas de base, la
fonction s'appelle indéfiniment : c'est la \textbf{récursion infinie}. Le
programme finit par planter (\emph{stack overflow}, débordement de pile). Vérifie
\textbf{toujours} qu'il existe un cas de base atteignable.
\end{attention}

\subsection{La factorielle}

Rappel mathématique : $n! = n \times (n-1) \times \dots \times 2 \times 1$, avec
$0! = 1$. On remarque que $n! = n \times (n-1)!$. Le cas de base est $0! = 1$.

\begin{codec}
// Factorielle de n, version recursive
int factorielle(int n) {
    if (n == 0) {            // CAS DE BASE
        return 1;
    }
    return n * factorielle(n - 1);   // CAS RECURSIF
}
\end{codec}

\subsection{La pile d'appels}

Pour calculer \code{factorielle(4)}, l'ordinateur empile les appels, puis les
résout en remontant :

\begin{exemple}
\code{factorielle(4)} = 4 $\times$ \code{factorielle(3)}\\
\hspace*{2em}= 4 $\times$ (3 $\times$ \code{factorielle(2)})\\
\hspace*{2em}= 4 $\times$ (3 $\times$ (2 $\times$ \code{factorielle(1)}))\\
\hspace*{2em}= 4 $\times$ (3 $\times$ (2 $\times$ (1 $\times$ \code{factorielle(0)})))\\
\hspace*{2em}= 4 $\times$ 3 $\times$ 2 $\times$ 1 $\times$ \textbf{1} = \textbf{24}
\end{exemple}

Chaque appel est mis en attente dans une zone mémoire appelée la \textbf{pile}
(\emph{stack}) tant que le suivant n'a pas renvoyé son résultat. Quand le cas de
base \code{factorielle(0)} renvoie \code{1}, la pile se vide en sens inverse.

\subsection{La puissance}

$x^n$ se calcule récursivement : $x^0 = 1$ (cas de base) et
$x^n = x \times x^{n-1}$.

\begin{codec}
// Calcule x puissance n (n entier positif), version recursive
int puissance(int x, int n) {
    if (n == 0) {            // CAS DE BASE : x^0 = 1
        return 1;
    }
    return x * puissance(x, n - 1);   // CAS RECURSIF
}
\end{codec}

\subsection{Somme de 1 à n}

La fonction \code{somme(n)} calcule $1 + 2 + \dots + n$. Cas de base :
\code{somme(0)} vaut \code{0} ; cas récursif : \code{somme(n) = n + somme(n-1)}.

\begin{codec}
// Somme des entiers de 1 a n, version recursive
int somme(int n) {
    if (n == 0) {
        return 0;            // CAS DE BASE
    }
    return n + somme(n - 1); // CAS RECURSIF
}
\end{codec}

\subsection{La suite de Fibonacci}

La suite de Fibonacci est définie par \code{F(0) = 0}, \code{F(1) = 1} et
\code{F(n) = F(n-1) + F(n-2)} pour \code{n >= 2}. Elle a donc \textbf{deux} cas
de base.

\begin{codec}
// n-ieme terme de la suite de Fibonacci, version recursive
int fibonacci(int n) {
    if (n == 0) return 0;    // CAS DE BASE
    if (n == 1) return 1;    // CAS DE BASE
    return fibonacci(n - 1) + fibonacci(n - 2);   // deux appels
}
\end{codec}

\begin{remarque}
La version récursive de Fibonacci est élégante mais \textbf{lente} pour les grands
\code{n} : elle recalcule plusieurs fois les mêmes termes. Pour \code{fibonacci(40)},
elle fait des millions d'appels. La version itérative (une boucle) est alors bien
plus efficace.
\end{remarque}

\subsection{Itératif ou récursif ?}

Tout calcul récursif peut s'écrire avec une boucle (\textbf{itératif}), et
réciproquement. Comparons la factorielle.

\begin{codec}
// Factorielle, version ITERATIVE (avec une boucle)
int factorielleIter(int n) {
    int produit = 1;
    for (int i = 2; i <= n; i++) {
        produit = produit * i;
    }
    return produit;
}
\end{codec}

\begin{center}\small
\begin{tabular}{@{}lll@{}}
\toprule
\textbf{Critère} & \textbf{Récursif} & \textbf{Itératif}\\
\midrule
Lisibilité & souvent plus proche des maths & plus terre à terre\\
Mémoire & consomme la pile (un cadre par appel) & quasi nulle\\
Vitesse & plus lent (appels) & en général plus rapide\\
Risque & récursion infinie / débordement & boucle infinie\\
\bottomrule
\end{tabular}
\end{center}

\begin{methode}
\textbf{Écrire une fonction récursive.} (1)~Identifier le \textbf{cas de base}
(la plus petite valeur, où la réponse est immédiate). (2)~Exprimer le résultat de
\code{f(n)} en fonction de \code{f(n-1)} (ou d'un cas plus petit). (3)~Vérifier
que chaque appel \textbf{rapproche} du cas de base.
\end{methode}

\section{Exercices}

\rubrique{Application}

\exo{Minimum de deux entiers}\diff{1}
Écris une fonction \code{int min(int a, int b)} qui renvoie le plus petit des deux
entiers. Teste-la dans \code{main} avec \code{min(8, 3)}.

\exo{Valeur absolue}\diff{1}
Écris une fonction \code{int valeurAbsolue(int n)} qui renvoie la valeur absolue de
\code{n} (par exemple \code{-7} donne \code{7}).

\exo{Mention au baccalauréat}\diff{2}
Écris une fonction \code{void afficherMention(float moyenne)} qui affiche la mention
obtenue : \emph{Passable} si la moyenne est entre 10 et 12, \emph{Assez bien} entre
12 et 14, \emph{Bien} entre 14 et 16, \emph{Très bien} à partir de 16, et
\emph{Recalé} en dessous de 10 (notes sur 20).

\exo{Factorielle dans les deux styles}\diff{2}
Écris la fonction \code{factorielle} en version \textbf{récursive}, puis en version
\textbf{itérative}. Appelle les deux dans \code{main} pour vérifier qu'elles
donnent le même résultat pour \code{5}.

\rubrique{Entraînement}

\exo{Nombre premier}\diff{2}
Écris une fonction \code{int estPremier(int n)} qui renvoie \code{1} si \code{n} est
un nombre premier, \code{0} sinon. Rappel : un nombre premier n'est divisible que
par \code{1} et par lui-même, et est strictement supérieur à \code{1}.

\exo{PGCD}\diff{3}
Écris une fonction \code{int pgcd(int a, int b)} qui calcule le plus grand commun
diviseur de \code{a} et \code{b} par l'algorithme d'Euclide (soustractions
successives, ou \code{pgcd(a, b) = pgcd(b, a \% b)} avec \code{pgcd(a, 0) = a}).

\exo{Somme des chiffres}\diff{3}
Écris une fonction \textbf{récursive} \code{int sommeChiffres(int n)} qui renvoie la
somme des chiffres de \code{n}. Par exemple \code{sommeChiffres(237)} vaut
\code{2 + 3 + 7 = 12}. Indice : le dernier chiffre est \code{n \% 10}, et le reste
du nombre est \code{n / 10}.

\exo{Puissance}\diff{2}
Écris une fonction \code{puissance(int x, int n)} en version récursive, puis appelle
\code{puissance(2, 10)} dans \code{main} (le résultat doit valoir \code{1024}).

\section{Corrigés}

\corrige{de l'exercice 1}
\begin{codec}
#include <stdio.h>

int min(int a, int b) {
    if (a < b) {
        return a;
    }
    return b;
}

int main(void) {
    printf("Le minimum est %d\n", min(8, 3));   // affiche 3
    return 0;
}
\end{codec}

\corrige{de l'exercice 2}
\begin{codec}
#include <stdio.h>

int valeurAbsolue(int n) {
    if (n < 0) {
        return -n;     // -(-7) = 7
    }
    return n;
}

int main(void) {
    printf("%d\n", valeurAbsolue(-7));   // affiche 7
    printf("%d\n", valeurAbsolue(4));    // affiche 4
    return 0;
}
\end{codec}

\corrige{de l'exercice 3}
\begin{codec}
#include <stdio.h>

void afficherMention(float moyenne) {
    if (moyenne < 10) {
        printf("Recale\n");
    } else if (moyenne < 12) {
        printf("Passable\n");
    } else if (moyenne < 14) {
        printf("Assez bien\n");
    } else if (moyenne < 16) {
        printf("Bien\n");
    } else {
        printf("Tres bien\n");
    }
}

int main(void) {
    afficherMention(13.5);   // affiche Assez bien
    afficherMention(16.2);   // affiche Tres bien
    return 0;
}
\end{codec}

\corrige{de l'exercice 4}
\begin{codec}
#include <stdio.h>

// Version recursive
int factorielle(int n) {
    if (n == 0) {
        return 1;
    }
    return n * factorielle(n - 1);
}

// Version iterative
int factorielleIter(int n) {
    int produit = 1;
    for (int i = 2; i <= n; i++) {
        produit = produit * i;
    }
    return produit;
}

int main(void) {
    printf("Recursive : %d\n", factorielle(5));      // 120
    printf("Iterative : %d\n", factorielleIter(5));  // 120
    return 0;
}
\end{codec}

\corrige{de l'exercice 5}
\begin{codec}
#include <stdio.h>

int estPremier(int n) {
    if (n < 2) {
        return 0;            // 0 et 1 ne sont pas premiers
    }
    for (int d = 2; d < n; d++) {
        if (n % d == 0) {    // n est divisible par d
            return 0;        // donc pas premier
        }
    }
    return 1;                // aucun diviseur trouve : premier
}

int main(void) {
    printf("%d\n", estPremier(7));    // 1 (premier)
    printf("%d\n", estPremier(12));   // 0 (non premier)
    return 0;
}
\end{codec}

\corrige{de l'exercice 6}
\begin{codec}
#include <stdio.h>

// Algorithme d'Euclide, version recursive
int pgcd(int a, int b) {
    if (b == 0) {
        return a;            // CAS DE BASE
    }
    return pgcd(b, a % b);   // CAS RECURSIF
}

int main(void) {
    printf("PGCD(48, 36) = %d\n", pgcd(48, 36));   // 12
    return 0;
}
\end{codec}

\corrige{de l'exercice 7}
\begin{codec}
#include <stdio.h>

// Somme des chiffres, version recursive
int sommeChiffres(int n) {
    if (n == 0) {
        return 0;                  // CAS DE BASE
    }
    return (n % 10) + sommeChiffres(n / 10);   // dernier chiffre + le reste
}

int main(void) {
    printf("%d\n", sommeChiffres(237));   // 12
    return 0;
}
\end{codec}

\corrige{de l'exercice 8}
\begin{codec}
#include <stdio.h>

int puissance(int x, int n) {
    if (n == 0) {
        return 1;
    }
    return x * puissance(x, n - 1);
}

int main(void) {
    printf("2^10 = %d\n", puissance(2, 10));   // 1024
    return 0;
}
\end{codec}

\begin{aretenir}
Une fonction se \textbf{déclare} (\code{type nom(params)\{ \dots return v; \}}),
s'\textbf{appelle} par son nom et reçoit ses arguments \textbf{par valeur} (une
copie). Place un \textbf{prototype} avant \code{main}. Une fonction
\textbf{récursive} s'appelle elle-même : elle a \textbf{toujours} un \textbf{cas
de base} pour s'arrêter, sinon c'est la récursion infinie.
\end{aretenir}
